% Auto-generated from Chapter-21-The-Refugee-Crisis-Report.md
% Source: C:\Users\PCGAME\Desktop\story&universes\chain-weapon-story\finished-manuscript\manuscriptbaseforchapters\Chapter-21-The-Refugee-Crisis-Report.md

\chapter{The Refugee Crisis Report}
\renewcommand{\CurrentChapterNum}{21}
\POV{\glyphAisen}{Aisen}

The woman was sitting against the granary wall when Aisen found her.

Not sitting—that word implied choice, the selection of a place and the decision to occupy it. What the woman was doing against the granary wall was something closer to arrival, the final collapse of forward momentum into stillness, the body reaching the end of its capacity for movement and simply stopping where it stopped, the way a stone thrown across water eventually stops skipping and sinks. She had a child in her arms. The child was asleep—or unconscious, Aisen could not tell from this distance, the difference between the two states visible only in the quality of the breathing, the depth of it, and the evening light was too poor for that kind of reading. The woman's head was tilted back against the rough stone, and her eyes were open but not seeing, fixed on a middle distance that contained nothing because it was not a place she was looking at but a place she was looking away from.

The town of Ashenmere sat three miles down the hill from the academy, connected by a road that the institution maintained just well enough to signal that the town existed within its sphere of influence and just poorly enough to signal that the town's needs were secondary to the academy's purposes. Students went there for market-day provisions, for the kind of gossip that did not circulate within institutional walls, for the tavern on the square that served beer so bad that drinking it constituted a minor rebellion against the academy's standards of excellence. The granary occupied the town's eastern edge, where the buildings thinned and the land began its slow concession to the scrublands that stretched toward the border. It was not a place that drew attention. That, Aisen understood with a clarity that felt like weight settling in his chest, was the point.

He had come to Ashenmere for ink—the blue-black variety that the academy's supply office stocked irregularly and that the stationer on the square kept in consistent supply, the kind of practical errand that gave a student reason to be off campus without requiring explanation. But the ink had been a secondary motive grafted onto a primary one. For three days, reports had been filtering through Rin's information channels—whispered fragments gathered from market traders, from a carter she had cultivated as a source, from the shifting patterns of road traffic that she read the way Kaelen read root systems—that people were moving through the border scrublands in numbers that the official reports did not account for. People moving at night. People moving without roads. People moving in the direction of anywhere that was not where they had been.

Aisen approached the woman. His footsteps on the packed dirt were deliberate—audible enough to announce himself, slow enough to not alarm. He had learned this calibration in the tavern years ago, the careful introduction of presence into the proximity of someone who had reason to fear the approach of strangers. The woman's eyes tracked the sound before they tracked him, her head turning a fraction of a second after her gaze, the sequence reversed from the way alert people moved—she was hearing the world before she was seeing it, functioning on the senses that required less effort, the body conserving what it could.

"There's a well behind the cooper's workshop," Aisen said. "Forty paces south. The water's clean."

The woman stared at him. Her face was thin in the way that faces become thin not from a day's hunger but from weeks of it, the fat reserves beneath the skin consumed so that the architecture of the skull became legible—cheekbones, jaw, the orbital ridges that usually live beneath the surface of a face but were now, in this woman's face, the dominant features, the structural truth that flesh normally conceals. She was perhaps thirty. Perhaps forty. The kind of exhaustion she carried made age a speculation rather than an observation. On her neck, below the left ear, a burn scar twisted the skin into a knot of tissue that was still pink at the edges—recent enough to be tender, old enough to have begun its permanent reshaping of the surface it occupied.

"The water's clean," Aisen repeated, because the woman had not responded, and because repetition was sometimes the bridge between hearing and comprehension when exhaustion had widened the gap between the two.

"I had water," she said. Her voice was hoarse and flat, pitched in the register of someone who had stopped modulating for expression and was simply conveying information at the minimum cost of effort. "At the river. Three days ago. I had water there."

Three days. Aisen let the number settle. The nearest river was the Selthe, which ran sixteen miles east through the scrublands before turning south toward the border. Sixteen miles in three days, carrying a child, without water since the crossing.

"Stay here," he said. "I'll be back."

He found Kaelen at the southeast edge of town, where the scrublands pushed against the last garden plot like a slow tide against a failing levee. Kaelen was crouched among the wild plants that grew in the margin between cultivated ground and waste ground—the interstitial space that interested him more than either extreme, the place where domesticated soil and wild soil conducted their negotiations. His fingers were in the earth. His green-gold eyes had the half-focused quality they took on during communion with root systems, the expression that Aisen had learned to read as concentration rather than absence.

"There's a woman at the granary," Aisen said. "With a child. She hasn't had water in three days."

Kaelen's hands lifted from the soil. The transition was immediate—the gentle focus of the nature-mage replaced by the sharper attention of someone being told a practical thing that required a practical response. Dirt fell from his fingers in small cascades. "Food?"

"I didn't ask. But from the look of her, longer than three days."

Kaelen stood. He looked at the wild margin where he'd been working, and Aisen watched the calculation happen—Kaelen's particular kind of calculation, which was not mathematical but ecological, an assessment of what the land held and what could be taken from it without damage. "There's burdock root along the Selthe path. Dock leaves. Wild garlic if you know where. I can gather enough for a meal in twenty minutes." He paused. "A hot meal needs fire. Fire needs a place that isn't visible."

"I know a place," Aisen said, though what he meant was: \emph{I know someone who knows places.}